\documentclass[a4paper,10pt]{article}
\usepackage[german]{babel} % for correct language and hyphenation and stuff
\usepackage{calc}           % for easier length calculations (infix notation)
\usepackage{enumitem}       % for configuring list environments
\usepackage{fancyhdr}       % for setting header and footer
\usepackage{fontspec}       % for fonts
\usepackage{geometry}       % for setting margins (\newgeometry)
\usepackage{graphicx}       % for pictures
\usepackage{lastpage}       % for getting total number of pages (LastPage)
\usepackage{microtype}      % for microtypography stuff
\usepackage{xcolor}         % for colours
\usepackage{newpxtext,newpxmath} % Palatine font

% margin and column widths
% ------------------------

% margins
\newgeometry{includehead,includefoot,left=15mm,right=15mm,top=15mm,bottom=15mm}

% width of the gap between left and right column
\newlength{\cvcolumngapwidth}
\setlength{\cvcolumngapwidth}{3.5mm}

% left column width
\newlength{\cvleftcolumnwidth}
\setlength{\cvleftcolumnwidth}{49mm}

% right column width
\newlength{\cvrightcolumnwidth}
\setlength{\cvrightcolumnwidth}{\textwidth-\cvleftcolumnwidth-\cvcolumngapwidth}

% set paragraph indentation to 0, because it screws up the whole layout otherwise
\setlength{\parindent}{0mm}


% style definitions
% -----------------
% style categories explanation:
% * \cvnameXXX is used for the name;
% * \cvsectionXXX is used for section names (left column, accompanied by a horizontal rule);
% * \cvtitleXXX is used for job/education titles (right column);
% * \cvdurationXXX is used for job/education durations (left column);
% * \cvheadingXXX is used for headings (left column);
% * \cvmainXXX (and \setmainfont) is used for main text;
% * \cvmarginXXX is used for text in margins;
% * \cvruleXXX is used for the horizontal rules denoting sections.

% font families
\defaultfontfeatures{Ligatures=TeX} % reportedly a good idea, see https://tex.stackexchange.com/a/37251

% colours
\definecolor{KITgreen}{RGB}{0,150,130}
\definecolor{KITblue}{RGB}{70,100,170}
\colorlet{cvnamecolor}{black}
\colorlet{cvsectioncolor}{KITgreen}
\colorlet{cvtitlecolor}{KITblue}
\colorlet{cvdurationcolor}{KITblue}
\colorlet{cvheadingcolor}{KITblue}
\colorlet{cvmaincolor}{black}
\colorlet{cvmargincolor}{black}
\colorlet{cvrulecolor}{KITgreen}

\color{cvmaincolor}

% styles
\newcommand{\cvnamestyle}[1]{{\Large\textcolor{cvnamecolor}{#1}}}
\newcommand{\cvsectionstyle}[1]{{\normalsize\textcolor{cvsectioncolor}{#1}}}
\newcommand{\cvtitlestyle}[1]{{\large\textcolor{cvtitlecolor}{#1}}}
\newcommand{\cvdurationstyle}[1]{{\small\textcolor{cvdurationcolor}{#1}}}
\newcommand{\cvheadingstyle}[1]{{\normalsize\textcolor{cvheadingcolor}{#1}}}


% inter-item spacing
% ------------------

% vertical space after personal info and standard CV items
\newlength{\cvafteritemskipamount}
\setlength{\cvafteritemskipamount}{5mm plus 1.25mm minus 1.25mm}

% vertical space after sections
\newlength{\cvaftersectionskipamount}
\setlength{\cvaftersectionskipamount}{2mm plus 0.5mm minus 0.5mm}

% extra vertical space to be used when a section starts with an item with a heading (e.g. in the skills section),
% so that the heading does not follow the section name too closely
\newlength{\cvbetweensectionandheadingextraskipamount}
\setlength{\cvbetweensectionandheadingextraskipamount}{1mm plus 0.25mm minus 0.25mm}


% intra-item spacing
% ------------------

% vertical space after name
\newlength{\cvafternameskipamount}
\setlength{\cvafternameskipamount}{3mm plus 0.75mm minus 0.75mm}

% vertical space after personal info lines
\newlength{\cvafterpersonalinfolineskipamount}
\setlength{\cvafterpersonalinfolineskipamount}{2mm plus 0.5mm minus 0.5mm}

% vertical space after titles
\newlength{\cvaftertitleskipamount}
\setlength{\cvaftertitleskipamount}{1mm plus 0.25mm minus 0.25mm}

% value to be used as parskip in right column of CV items and itemsep in lists (same for both, for consistency)
\newlength{\cvparskip}
\setlength{\cvparskip}{0.5mm plus 0.125mm minus 0.125mm}

% set global list configuration (use parskip as itemsep, and no separation otherwise)
\setlist{parsep=0mm,topsep=0mm,partopsep=0mm,itemsep=\cvparskip}


% CV commands
% -----------

% creates a "personal info" CV item with the given left and right column contents, with appropriate vertical space after
% @param #1 left column content (should be the CV photo)
% @param #2 right column content (should be the name and personal info)
\newcommand{\cvpersonalinfo}[2]{
    % left and right column
    \begin{minipage}[t]{\cvleftcolumnwidth}
        \vspace{0mm} % XXX hack to align to top, see https://tex.stackexchange.com/a/11632
        \raggedleft #1
    \end{minipage}% XXX necessary comment to avoid unwanted space
    \hspace{\cvcolumngapwidth}% XXX necessary comment to avoid unwanted space
    \begin{minipage}[t]{\cvrightcolumnwidth}
        \vspace{0mm} % XXX hack to align to top, see https://tex.stackexchange.com/a/11632
        #2
    \end{minipage}

    % space after
    \vspace{\cvafteritemskipamount}
}

% typesets a name, with appropriate vertical space after
% @param #1 name text
\newcommand{\cvname}[1]{
    % name
    \cvnamestyle{#1}

    % space after
    \vspace{\cvafternameskipamount}
}

% typesets a line of personal info beginning with an icon, with appropriate vertical space after
% @param #1 parameters for the \includegraphics command used to include the icon
% @param #2 icon filename
% @param #3 line text
\newcommand{\cvpersonalinfolinewithicon}[3]{
    % icon, vertically aligned with text (see https://tex.stackexchange.com/a/129463)
    \raisebox{.5\fontcharht\font`E-.5\height}{\includegraphics[#1]{#2}}
    % text
    #3

    % space after
    \vspace{\cvafterpersonalinfolineskipamount}
}

% creates a "section" CV item with the given left column content, a horizontal rule in the right column, and with
% appropriate vertical space after
% @param #1 left column content (should be the section name)
\newcommand{\cvsection}[1]{
    % left and right column
    \begin{minipage}[t][][b]{\cvleftcolumnwidth}
        \raggedleft\cvsectionstyle{#1}
    \end{minipage}% XXX necessary comment to avoid unwanted space
    \hspace{\cvcolumngapwidth}% XXX necessary comment to avoid unwanted space
    \begin{minipage}[t]{\cvrightcolumnwidth}
        \textcolor{cvrulecolor}{\rule{\cvrightcolumnwidth}{0.3mm}}
    \end{minipage}

    % space after
    \vspace{\cvaftersectionskipamount}
}

% creates a standard, multi-purpose CV item with the given left and right column contents, parskip set to cvparskip
% in the right column, and with appropriate vertical space after
% @param #1 left column content
% @param #2 right column content
\newcommand{\cvitem}[2]{
    % left and right column
    \begin{minipage}[t]{\cvleftcolumnwidth}
    \strut\vspace*{-\baselineskip}\newline %quick fix to align both minipages to top baseline
    \raggedleft #1
    \end{minipage}% XXX necessary comment to avoid unwanted space
    \hspace{\cvcolumngapwidth}% XXX necessary comment to avoid unwanted space
    \begin{minipage}[t]{\cvrightcolumnwidth}
        \setlength{\parskip}{\cvparskip}
        \strut\vspace*{-\baselineskip}\newline #2 %quick fix to align both minipages to top baseline
    \end{minipage}

    % space after
    \vspace{\cvafteritemskipamount}
}

% typesets a title, with appropriate vertical space after
% @param #1 title text
\newcommand{\cvtitle}[1]{
    % title
    \cvtitlestyle{#1}

    % space after
    \vspace{\cvaftertitleskipamount}
    % XXX need to subtract cvparskip here, because it is automatically inserted after the title "paragraph"
    \vspace{-\cvparskip}
}


% header and footer
% -----------------

% header
\renewcommand{\headrulewidth}{0mm} % needed to remove line under header that is there by default
\fancypagestyle{firstpage}{
    \fancyhead[R]{} % no name in header on the first page
    \fancyhead[L]{\textcolor{cvmargincolor}{Curriculum vit\ae}}
}
\fancypagestyle{otherpages}{
    \fancyhead[R]{\textcolor{cvmargincolor}{Olena Manzhura}}
    \fancyhead[L]{\textcolor{cvmargincolor}{Curriculum vit\ae}}
}
\thispagestyle{firstpage}
\pagestyle{otherpages}

% footer
\pagestyle{fancy}
\fancyfoot{} % needed to remove page number that is in centre of footer by default
\fancyfoot[L]{\footnotesize\textcolor{cvmargincolor}{\today}}
\fancyfoot[R]{\footnotesize\textcolor{cvmargincolor}{\thepage~/~\pageref{LastPage}}}

\begin{document}
\cvpersonalinfo{
    % text to fill space instead of a photo
    \cvsectionstyle{PERSÖNLICHE DATEN}
}{
    % name
    \cvname{Olena Manzhura}

    % address
    \cvpersonalinfolinewithicon{height=4mm}{resources/europasscv-icons/address_europass_icon.pdf}{
        Hardtstr. 22, 76185 Karlsruhe
    }

    % phone number
    \cvpersonalinfolinewithicon{height=4mm}{resources/europasscv-icons/mobile_europass_icon.pdf}{
        +49\;176\;232\;041\;82
    }

    % email address
    \cvpersonalinfolinewithicon{height=4mm}{resources/europasscv-icons/mail_europass_icon.pdf}{
        manzhuraolena@gmail.com
    }

    % personal website
    %\cvpersonalinfolinewithicon{height=4mm}{resources/europasscv-icons/website_europass_icon.pdf}{
    %    www.fake-name-homepage.com
    %}

    % date of birth
    Geboren am 08. November 1997 in Garkushynci, Ukraine
}

\cvsection{SCHULE UND STUDIUM}
\cvitem{\cvdurationstyle{2018 -- 2021}}{\cvtitle{M. Sc. Elektrotechnik und Informationstechnik}
	Karlsruher Institut für Technologie
	\begin{itemize}[leftmargin=*]
		\item Vertiefungsrichtung: Adaptronik
		\item Abschlussarbeit \textit{A Terabit Sampling System with a Photonics Time-Stretch ADC} (Note: 1.3)
		\item Abschlussnote: 1.8
	\end{itemize}
}

\cvitem{\cvdurationstyle{Januar 2020 -- Juni 2020}}{\cvtitle{Auslandssemester}
	Université de Technologie de Troyes, Frankreich \\
}

\cvitem{\cvdurationstyle{2015 -- 2018}}{\cvtitle{B. Sc. Elektrotechnik und Informationstechnik}
	Karlsruher Institut für Technologie
	\begin{itemize}[leftmargin=*]
		\item Abschlussarbeit: \textit{Kalibrierung eines Ganzkörperzählers mit HPGe Detektoren	zur Bestimmung der Aktivität im Körper von Menschen} (Note: 1.3)
		\item Abschlussnote: 2.1
	\end{itemize}
}

\cvitem{\cvdurationstyle{2007 -- 2015}}{\cvtitle{Gymnasium}
	Markgrafen-Gymnasium, Karlsruhe-Durlach
	\begin{itemize}[leftmargin=*]
		\item Abschlussnote: 1.8
	\end{itemize}
}

\cvsection{BERUFLICHE ERFAHRUNG}

\cvitem{\cvdurationstyle{April 2022 -- heute}}{\cvtitle{Wissenschaftliche Mitarbeiterin}
	Promotion am Institut für Prozessdatenverarbeitung und Elektronik (IPE) am KIT
	\begin{itemize}[leftmargin=*]
		\item Vorläufiger Titel der Arbeit: \textit{Real-Time Data Acquisition System for Terascale Era Experiments}
	\end{itemize}
}

\cvitem{\cvdurationstyle{September 2021 -- März 2022}}{\cvtitle{Wissenschaftliche Hilfskraft}
	Institut für Prozessdatenverarbeitung und Elektronik (IPE) am KIT
	\begin{itemize}[leftmargin=*]
		\item Entwicklung von Firmware auf einem FPGA für ein Terabit-Abtastsystem
	\end{itemize}
}

\cvitem{\cvdurationstyle{Mai 2017 -- Juni 2021}}{\cvtitle{Wissenschaftliche Hilfskraft}
	Institut für Elektroenergiesysteme und Hochspannungstechnik am KIT
	\begin{itemize}[leftmargin=*]
		\item Aufbau und Wartung einer Wetterstation
		\item Implementierung und Bewertung von Smart-Home-Lösungen für den Raspberry Pi
	\end{itemize}
}

\cvitem{\cvdurationstyle{Oktober 2018 -- Dezember 2019}}{\cvtitle{Wissenschaftliche Hilfskraft}
	Sicherheit und Umwelt (SUM) am KIT
	\begin{itemize}[leftmargin=*]
		\item Kalibrierung von Detektoren für In-vivo-Messungen von radioaktiven Stoffen
	\end{itemize}
}

\cvitem{\cvdurationstyle{April 2018 -- Juli 2018}}{\cvtitle{Praktikum}
	IWK Verpackungstechnik GmbH, Stutensee
	\begin{itemize}[leftmargin=*]
		\item Optimierung der Bewegungsbahnen eines Bestückungsroboters
		\item SPS-Programmierung
	\end{itemize}
}

\cvitem{\cvdurationstyle{April 2018 -- Juli 2018}}{\cvtitle{Tutor}
	Institut für Regelungs- und Steuerungssysteme (IRS) am KIT
	\begin{itemize}[leftmargin=*]
		\item Halten des Tutoriums für die Vorlesung \glqq Systemdynamik und Regelungstechnik\grqq
	\end{itemize}
}

\cvsection{KENNTNISSE}
%\vspace{\cvbetweensectionandheadingextraskipamount}

\cvitem{\cvheadingstyle{Sprachen}}{
	\hspace*{-2mm}\begin{tabular}{lcl}
		Deutsch  & -- & Muttersprache \\
		Russisch & -- & Muttersprache \\
		Englisch & -- & CEFR: C1, Cambridge Certificate of English \\
		Französisch & -- & CEFR: B2, Französischkurs an der UTT (Troyes, Frankreich) \\
	\end{tabular}
}

\cvitem{\cvheadingstyle{EDV}}{
    \\[-5.5ex]\begin{itemize}[leftmargin=*]
    	\item Programmieren: Python
    %	\item CAD: Altium Designer, PADS
   % 	\item HDL: Vivado, Verilog, VHDL
    	\item \LaTeX, MS Office
    	\item Grundlegende Linux-Administration
    \end{itemize}
}


\cvsection{WEITERE INFORMATIONEN}\vspace{\cvbetweensectionandheadingextraskipamount}
\cvitem{\cvheadingstyle{Führerschein}}{
	\\[-5.5ex]\begin{itemize}[leftmargin=*]
     \item Klasse B
    \end{itemize}
}

\cvitem{\cvheadingstyle{Ehrenamt}}{
	\\[-5.5ex]\begin{itemize}[leftmargin=*]
	        \item Technisches Hilfswerk, OV Karlsruhe: Fachgruppe Räumen
	        \item DLRG OG Durlach: Jugendvorstand, Rettungsschwimmer
	        \item Freiwillige Feuerwehr Karlsruhe Abteilung Mühlburg: Frauenbeauftragte, Jugendwartin
	\end{itemize}
}

\end{document}
